\documentclass[12pt]{article}
\usepackage[utf8]{inputenc}
\usepackage[T1]{fontenc}
\usepackage{hyperref}
\title{Lovono Language Resources for K-12 Curriculum Development}
\author{Compiled by Claude Sonnet 5 (claude-sonnet-5) for Andr\'as Kornai}
\date{}
\begin{document}
\maketitle

\section{Introduction}
Lovono (ISO 639-3 vnk; also Vano, Whanou) is an Austronesian language of Vanikoro, Solomon Islands. Sources give 4 speakers (2012) or 5, and UNESCO rates it critically endangered. Teanu, the island's dominant language, has largely replaced it. Documentation is mostly linguistic.

\section{Primary Educational Resources}
No educational materials or school programmes were found.

\section{Linguistic and Reference Materials}
\begin{itemize}
\item Wikipedia article on Lovono, which mentions the grammar (2009) and lexicon (2021) work of A. Fran\c{c}ois on the Vanikoro languages; these works were named in the search summary but not located directly: \url{https://en.wikipedia.org/wiki/Lovono_language}
\item Wiktionary category for Lovono: \url{https://en.wiktionary.org/wiki/Category:Lovono_language}
\item Verbix page for Lovono: \url{https://docs.verbix.com/Languages/Lovono}
\item Fran\c{c}ois fieldwork and languages page: \url{https://marama.huma-num.fr/AF-field.htm}
\end{itemize}

\section{Religious and Community Resources}
None were found.

\section{Digital and Audio Resources}
\begin{itemize}
\item PARADISEC item DT1-112, a recorded 400-word list of Vano from Vanikoro collected by Darrell Tryon: \url{https://catalog.nabu-prod.paradisec.org.au/collections/DT1/items/112}
\item Global Recordings Network page for the language: \url{https://globalrecordings.net/en/language/vnk}
\item Wikitongues page for vnk: \url{https://wikitongues.org/languages/vnk/}
\end{itemize}

\section{Curriculum Development Context}
With about four elderly speakers, a school curriculum is not feasible in the usual sense. Documentation and archiving are the realistic goals: the word-list recording and the linguistic literature are the main sources to work from. Any revitalization effort would need to be led by the community and by linguists with the relevant fieldwork.
\end{document}